\chapter{Manual de instalaci\'on}\label{apendice-instalacion}
% tfg-floatbarrier-auto
\makeatletter
\@ifundefined{tfgfloatbarrier}{
  \def\tfgfloatbarrier{1}
  \let\tfgorigfigure\figure
  \let\tfgendorigfigure\endfigure
  \renewenvironment{figure}{\tfgorigfigure}{\tfgendorigfigure\FloatBarrier}
  \let\tfgorigtable\table
  \let\tfgendorigtable\endtable
  \renewenvironment{table}{\tfgorigtable}{\tfgendorigtable\FloatBarrier}
}{}
\makeatother

\section{Requisitos previos}

Para ejecutar la soluci\'on en entorno local se requiere JDK 21 y Maven 3.9+ para backend, Node.js 18+ y npm 9+ para frontend, y PostgreSQL 14+ disponible en local o contenedor. Estos requisitos aseguran compatibilidad con el stack desplegado.

\section{Preparaci\'on del entorno}

La preparaci\'on del entorno consiste en clonar el repositorio y situarse en la ra\'iz del proyecto, y en verificar credenciales y variables seg\'un el Anexo B. Este orden evita configuraciones incompletas antes de arrancar servicios.

\section{Base de datos}

Para entorno local se recomienda PostgreSQL con la base \texttt{tfgdb}
y el usuario \texttt{tfg}. El repositorio incluye un script de apoyo
\texttt{backend/run-postgres.ps1} que comprueba el servicio y arranca
el backend. En caso de preferir contenedor, puede utilizarse
\texttt{docker-compose} (ver subsecci\'on correspondiente).

\section{Arranque del backend}

Pasos de ejecuci\'on recomendados:

\begin{verbatim}
cd backend
./mvnw clean install
./mvnw spring-boot:run
\end{verbatim}

Si se desea forzar configuraci\'on local con PostgreSQL,
puede activarse el perfil \texttt{local} mediante
\texttt{SPRING\_PROFILES\_ACTIVE=local}. Para pruebas r\'apidas
sin base externa, est\'a disponible el perfil \texttt{h2}.

\section{Arranque del frontend}

\begin{verbatim}
cd frontend
npm install
npm run dev
\end{verbatim}

Por defecto el frontend consumir\'a la API local. Si es necesario,
ajustar \texttt{VITE\_API\_BASE\_URL} seg\'un el entorno.

\section{Ejecuci\'on con contenedores}

Como alternativa reproducible, la soluci\'on puede ejecutarse con
\texttt{docker-compose}, levantando base de datos y backend con healthchecks.
Esta opci\'on reduce diferencias entre m\'aquinas de desarrollo.

\section{Build de producci\'on}

Para generar artefactos de despliegue:

\begin{verbatim}
cd backend
./mvnw clean package -DskipTests

cd frontend
npm run build
\end{verbatim}

\FloatBarrier
\chapter{Configuraci\'on por entornos}\label{apendice-configuracion}

\section{Variables cr\'iticas de backend}

Las variables m\'inimas para el funcionamiento seguro en entornos no locales incluyen los par\'ametros de conexi\'on a la base de datos (\texttt{DATABASE\_URL}, \texttt{DATABASE\_USER}, \texttt{DATABASE\_PASSWORD}), la clave privada para la firma criptogr\'afica de tokens de sesi\'on (\texttt{JWT\_SECRET}), el perfil de entorno activo para la inyecci\'on de dependencias (\texttt{SPRING\_PROFILES\_ACTIVE}), y la URL base del cliente web para configurar correctamente las reglas CORS (\texttt{APP\_FRONTEND\_BASE\_URL}). Esta configuraci\'on garantiza conectividad, seguridad y coherencia de entorno.

En producci\'on se habilitan adem\'as variables para servicios externos
como almacenamiento cloud y correo transaccional. En el despliegue sobre
Render, el correo se configura mediante Brevo usando API HTTP, por lo que
deben definirse \texttt{BREVO\_API\_KEY} y \texttt{BREVO\_FROM\_EMAIL}
con un remitente verificado en Brevo. Esta elecci\'on evita depender de
conexiones SMTP salientes, que pueden estar restringidas por el PaaS.

\section{Variables cr\'iticas de frontend}

La variable principal de configuraci\'on de cliente es
\texttt{VITE\_API\_BASE\_URL}, que define el endpoint de API consumido.

\FloatBarrier
\chapter{Operaci\'on y verificaci\'on}\label{apendice-operacion}

\section{Comprobaciones de salud}

El sistema expone endpoints de salud para diferenciaci\'on operativa: \'el de \textit{liveness} valida disponibilidad del proceso y el de \textit{readiness} valida disponibilidad real de dependencias como la base de datos. Esta separaci\'on permite detectar fallos de conectividad de forma temprana.

Esta separaci\'on permite detectar de forma temprana incidencias de conectividad
y reducir falsos positivos en monitorizaci\'on.

\section{Validaci\'on post-despliegue}

Tras cada despliegue se recomienda, como m\'inimo, comprobar endpoints de salud, realizar login con usuario de prueba, ejecutar operaci\'on de lectura y escritura en flujo acad\'emico, verificar subida y descarga de material y probar el correo transaccional (recuperaci\'on de contrase\~na o notificaci\'on) mediante Brevo. Estas comprobaciones validan la operativa esencial en producci\'on.

\FloatBarrier
\chapter{Declaraci\'on de uso de IA generativa}\label{apendice-ia}

En este TFG se ha utilizado IA generativa como herramienta de apoyo
en tareas acotadas y verificables, con el objetivo de mejorar claridad
documental, revisar estructura y acelerar actividades de bajo riesgo.
La IA no ha sustituido el criterio t\'ecnico ni ha generado decisiones
de dise\~no o implementaci\'on sin validaci\'on humana.

\section{Alcance y criterios de uso}

El uso de IA se limit\'o a apoyo de redacci\'on y reformulaci\'on de texto t\'ecnico, revisi\'on de estructura y coherencia de secciones, generaci\'on de listados de verificaci\'on y casos borde, microcopy y textos de ayuda para interfaz, y aclaraciones puntuales de sintaxis o formatos (por ejemplo, JSON o \LaTeX{}). Estas tareas se consideraron de bajo riesgo y siempre se validaron manualmente.

Se excluyeron expl\'icitamente usos como decisiones arquitect\'onicas finales sin justificaci\'on propia, c\'odigo productivo incorporado sin revisi\'on manual, resultados emp\'iricos o m\'etricas no verificadas y contenido no contrastado con el repositorio y la evidencia de pruebas. Esta restricci\'on garantiza responsabilidad t\'ecnica.

\section{Proceso de control y verificaci\'on}

Cada propuesta generada por IA se trat\'o como borrador. El proceso seguido consisti\'o en formular la consulta indicando el contexto del proyecto, evaluar si la respuesta era coherente con requisitos y c\'odigo, ajustar manualmente terminolog\'ia y rigor acad\'emico, contrastar con artefactos del repositorio (memoria, c\'odigo, tests) e integrar la versi\'on final ya revisada. Este flujo evita incorporaciones autom\'aticas sin validaci\'on.

\section{Ejemplos concretos de consultas y respuestas}

Los siguientes ejemplos resumen consultas reales y la respuesta resumida
que se utiliz\'o como punto de partida. En todos los casos se realiz\'o
revisi\'on manual antes de incorporar el contenido a la memoria o al producto.

\textbf{Ejemplo 1. Redacci\'on acad\'emica de objetivos.}
\textbf{Consulta}: \emph{Reescribe este p\'arrafo de objetivos para un tono acad\'emico,
evitando repeticiones y manteniendo el significado}. 
\textbf{Respuesta resumida}: propuesta de redacci\'on m\'as concisa
con estructura en dos frases y vocabulario t\'ecnico coherente.
\textbf{Aplicaci\'on}: se ajustaron t\'erminos a la nomenclatura del ERS
y se incorpor\'o solo tras validar la trazabilidad con requisitos.

\textbf{Ejemplo 2. Estructura de secci\'on comparativa.}
\textbf{Consulta}: \emph{Sugiere un esquema para un an\'alisis de competidores
en una memoria de TFG de software}. 
\textbf{Respuesta resumida}: se propuso dividir por tipolog\'ias de competidor,
incluir criterios de comparaci\'on y cerrar con diferenciadores propios.
\textbf{Aplicaci\'on}: el esquema se adapt\'o al contenido real del proyecto
y a las evidencias t\'ecnicas implementadas.

\textbf{Ejemplo 3. Validaci\'on de estructuras ZIP.}
\textbf{Consulta}: \emph{Prop\'on ejemplos de reglas para validar estructura
interna de un ZIP con archivos obligatorios y carpetas esperadas}. 
\textbf{Respuesta resumida}: lista de casos con archivos obligatorios,
extensiones permitidas y modos estricto/m\'inimo requerido.
\textbf{Aplicaci\'on}: se contrast\'o con la l\'ogica del servicio de validaci\'on
y se redact\'o la explicaci\'on en la memoria con terminolog\'ia consistente.

\textbf{Ejemplo 4. Textos de ayuda en interfaz de entrega.}
\textbf{Consulta}: \emph{Sugiere mensajes breves y claros para indicar
que una entrega requiere un ZIP con nombre y estructura concretos}. 
\textbf{Respuesta resumida}: propuestas de microcopy con tono informativo
y referencias expl\'icitas a nombre esperado, tama\~no y estructura.
\textbf{Aplicaci\'on}: los textos se ajustaron manualmente para alinearlos
con el vocabulario usado en la UI del proyecto.

\textbf{Ejemplo 5. Casos borde para pruebas.}
\textbf{Consulta}: \emph{Enumera casos borde para pruebas de subida de ZIP
en un sistema de entregables}. 
\textbf{Respuesta resumida}: listado de escenarios (nombre incorrecto,
estructura incompleta, archivos extra, ZIP corrupto, tama\~no excedido).
\textbf{Aplicaci\'on}: se utiliz\'o como checklist de verificaci\'on y
se descartaron casos no aplicables al alcance real.

Este anexo se incluye para garantizar transparencia metodol\'ogica
y trazabilidad del proceso de elaboraci\'on de la memoria.

\FloatBarrier
\chapter{Documentaci\'on complementaria disponible}\label{apendice-docs}

Como soporte adicional al contenido principal de la memoria, el repositorio oficial del proyecto incorpora documentos t\'ecnicos de operaci\'on, seguridad, testing y despliegue que permiten ampliar detalle en revisiones futuras. 

El c\'odigo fuente del proyecto se encuentra alojado en GitHub y est\'a disponible en el siguiente enlace: \url{https://github.com/Mancalrod/TFG}. 

Dentro del repositorio, se puede consultar la documentaci\'on t\'ecnica complementaria en la ruta \texttt{docs/}, donde se encuentran esquemas, configuraciones de despliegue y otros artefactos de dise\~no adicionales no incluidos expl\'icitamente en esta memoria.

\FloatBarrier
\chapter{Caso de uso de ejemplo}\label{apendice-caso-uso}

Para ilustrar de forma unificada la funcionalidad de la plataforma, se describe a continuaci\'on el recorrido completo (flujo de extremo a extremo) de una pr\'actica a trav\'es de los tres perfiles del sistema: Administrador, Profesor y Estudiante. Este recorrido es representativo de la operativa del TFG y coincide con el hilo conductor de la demostraci\'on en v\'ideo del proyecto.

\section{Administrador: Configuraci\'on inicial}

El administrador accede a la plataforma desde el enlace oficial \url{https://tfg-85os.onrender.com/admin} para realizar la configuraci\'on inicial del censo acad\'emico.

En primer lugar, crea el curso ``Ingenier\'ia de Software'' y le asigna un profesor responsable para su impartici\'on (Figura~\ref{fig:apendice-crear-curso}).

\figura{0.90}{\detokenize{../Pantallas de TFG/Crear curso apendice F.PNG}}{Creaci\'on del curso ``Ingenier\'ia de Software'' y asignaci\'on del profesorado.}{fig:apendice-crear-curso}{}

A continuaci\'on, procede a crear un grupo denominado ``Grupo A'', al cual le asocia la asignatura ``Ingenier\'ia de Software'' previamente establecida (Figura~\ref{fig:apendice-crear-grupo}).

\figura{0.90}{\detokenize{../Pantallas de TFG/Crear grupo y asignar curso apendice F.PNG}}{Creaci\'on de ``Grupo A'' y vinculaci\'on del curso de Ingenier\'ia.}{fig:apendice-crear-grupo}{}

Por \'ultimo, realiza la asignaci\'on de los estudiantes matriculados en dicho grupo (Figura~\ref{fig:apendice-anadir-estudiantes}). A estos alumnos los puede incorporar de manera interactiva seleccion\'andolos de la lista o bien cargando un archivo de texto plano (\texttt{.txt}) que contenga la lista de correos electr\'onicos correspondientes. Es un requisito indispensable del sistema que los estudiantes hayan sido previamente registrados como usuarios generales de la plataforma antes de poder ser elegidos o asignados a un grupo.

\figura{0.90}{\detokenize{../Pantallas de TFG/Añadir estudiante apendice F.PNG}}{Asignaci\'on e importaci\'on de estudiantes al ``Grupo A''.}{fig:apendice-anadir-estudiantes}{}

\section{Profesor: Creaci\'on de actividad y entregables}

El profesor inicia sesi\'on en el portal acad\'emico mediante la URL \url{https://tfg-85os.onrender.com/dashboard} y visualiza la asignatura ``Ingenier\'ia de Software'' asignada en su panel (Figura~\ref{fig:apendice-pantalla-inicial-profesor}).

\figura{0.90}{\detokenize{../Pantallas de TFG/pantalla inicial apendice f.PNG}}{Vista del panel principal docente con las asignaturas asignadas.}{fig:apendice-pantalla-inicial-profesor}{}

Tras entrar en la asignatura a trav\'es del enlace \url{https://tfg-85os.onrender.com/cursos/1}, crea una nueva \textbf{Actividad} (por ejemplo, ``Pr\'actica 1''), y la vincula estrictamente al ``Grupo A'' para garantizar que s\'olo los estudiantes matriculados en ese grupo tengan visibilidad y acceso a ella (Figura~\ref{fig:apendice-crear-actividad}).

\figura{0.90}{\detokenize{../Pantallas de TFG/Crear actividad apendice F.PNG}}{Formulario de creaci\'on de actividad y asignaci\'on de grupo.}{fig:apendice-crear-actividad}{}

Una vez dentro de los detalles de la actividad mediante el enlace \url{https://tfg-85os.onrender.com/actividades/1/entregables/nuevo}, define un \textbf{Entregable} concreto. En este formulario configura tanto la fecha l\'imite de entrega como las restricciones de formato esperadas, exigiendo de forma obligatoria que el env\'io sea un archivo comprimido en formato ZIP que contenga en su interior un documento PDF con la resoluci\'on (Figura~\ref{fig:apendice-crear-entregable-1} y Figura~\ref{fig:apendice-crear-entregable-2}).

\figura{0.90}{\detokenize{../Pantallas de TFG/Dentro de una actividad - Crear entregables (Parte 1) (Profesor).png}}{Configuraci\'on del entregable y fecha l\'imite (1/2).}{fig:apendice-crear-entregable-1}{}

\figura{0.90}{\detokenize{../Pantallas de TFG/Dentro de una actividad - Crear entregables (Parte 2) (Profesor).png}}{Configuraci\'on de restricciones de archivo obligatorias del entregable (2/2).}{fig:apendice-crear-entregable-2}{}

\section{Estudiante: Realizaci\'on de la entrega}

El estudiante inicia sesi\'on en la plataforma desde el enlace \url{https://tfg-85os.onrender.com/dashboard}. En su panel de control observa de inmediato que tiene una tarea pendiente dentro de la asignatura ``Ingenier\'ia de Software'' (Figura~\ref{fig:apendice-dashboard-alumno}).

\figura{0.90}{\detokenize{../Pantallas de TFG/Pantalla inicial alumno apendice F.PNG}}{Dashboard del estudiante con alertas de entregas pendientes.}{fig:apendice-dashboard-alumno}{}

Al acceder a los detalles del entregable mediante la URL \url{https://tfg-85os.onrender.com/entregables/13}, comprueba rigurosamente los requisitos impuestos por el docente: la fecha l\'imite y la estructura del ZIP requerida (Figura~\ref{fig:apendice-detalles-entregable-alumno}).

\figura{0.90}{\detokenize{../Pantallas de TFG/Activiad alumno apendice f.PNG}}{Detalles del entregable y estructura de validaci\'on del archivo ZIP.}{fig:apendice-detalles-entregable-alumno}{}

Tras preparar su archivo comprimido en su equipo local, el estudiante accede al formulario de env\'io a trav\'es de la ruta \url{https://tfg-85os.onrender.com/entregables/13/entregar}. Sube su archivo ZIP y pulsa enviar. El sistema realiza una validaci\'on sint\'actica y estructural en tiempo real (comprobando la extensi\'on y la presencia del PDF dentro del ZIP) y, tras superarla con \'exito, registra la entrega de forma segura en la base de datos y le muestra una confirmaci\'on visual (Figura~\ref{fig:apendice-realizar-entrega}).

\figura{0.90}{\detokenize{../Pantallas de TFG/Realizar entrega apendice f.PNG}}{Formulario de subida de archivos y validaci\'on en tiempo real.}{fig:apendice-realizar-entrega}{}

\section{Profesor: Correcci\'on y retroalimentaci\'on}

Transcurrido el plazo, el docente accede a la secci\'on global de evaluaciones desde el enlace \url{https://tfg-85os.onrender.com/evaluaciones} o ingresa directamente al detalle de la actividad para comprobar las entregas pendientes de correcci\'on (Figura~\ref{fig:apendice-evaluaciones-profesor}).

\figura{0.90}{\detokenize{../Pantallas de TFG/Correcion Entregable apendice f.PNG}}{Bandeja de entregas pendientes de calificaci\'on para el profesorado.}{fig:apendice-evaluaciones-profesor}{}

A continuaci\'on, ingresa a la entrega del estudiante desde la URL \url{https://tfg-85os.onrender.com/entregas/14}. En esta pantalla puede visualizar de forma integrada los metadatos de la entrega y el historial de versiones del archivo (Figura~\ref{fig:apendice-ver-entrega-profesor}), as\'i como previsualizar de forma interactiva el c\'odigo o contenido del archivo ZIP directamente en el navegador sin necesidad de descargarlo en local (Figura~\ref{fig:apendice-previsualizar-zip}).

\figura{0.90}{\detokenize{../Pantallas de TFG/visualizar entrega alumno apendice f.PNG}}{Vista detallada del material entregado e historial de env\'ios.}{fig:apendice-ver-entrega-profesor}{}

\figura{0.90}{\detokenize{../Pantallas de TFG/Previsalizar zip apendice f.PNG}}{Previsualizador integrado de la estructura interna del archivo ZIP.}{fig:apendice-previsualizar-zip}{}

Una vez analizado el trabajo, el profesor le asigna una calificaci\'on de ``8.5/10'' y redacta una evaluaci\'on detallada con observaciones de mejora en el cuadro de feedback (Figura~\ref{fig:apendice-calificar-trabajo}).

\figura{0.90}{\detokenize{../Pantallas de TFG/Calificar trabajo apendice f.PNG}}{Formulario de emisi\'on de notas y retroalimentaci\'on formativa.}{fig:apendice-calificar-trabajo}{}

Tras guardar y publicar la nota en el sistema mediante el enlace \url{https://tfg-85os.onrender.com/entregables/13}, las calificaciones pasan a estar visibles para los estudiantes de la asignatura (Figura~\ref{fig:apendice-mostrar-notas}).

\figura{0.90}{\detokenize{../Pantallas de TFG/Mostrar notas apendice f.PNG}}{Publicaci\'on de notas y cierre de evaluaci\'on del entregable.}{fig:apendice-mostrar-notas}{}

\section{Estudiante: Consulta de resultados}

El estudiante inicia sesi\'on de nuevo en la plataforma y, en su dashboard accesible desde \url{https://tfg-85os.onrender.com/dashboard}, recibe la notificaci\'on de que su nota ha sido publicada y est\'a disponible para su consulta (Figura~\ref{fig:apendice-nota-publicada}).

\figura{0.90}{\detokenize{../Pantallas de TFG/nota publicada apendice f.PNG}}{Dashboard del alumno con indicaci\'on de actividad calificada.}{fig:apendice-nota-publicada}{}

Al entrar al detalle de su entrega mediante el enlace \url{https://tfg-85os.onrender.com/entregas/14}, el estudiante visualiza con total transparencia y trazabilidad su calificaci\'on de ``8.5'' y el comentario explicativo aportado por el docente, concluyendo con \'exito el ciclo de evaluaci\'on (Figura~\ref{fig:apendice-visualizar-nota-alumno}).

\figura{0.90}{\detokenize{../Pantallas de TFG/Visualizar nota estudiante apendice f.PNG}}{Vista del estudiante con nota y comentarios de feedback del profesor.}{fig:apendice-visualizar-nota-alumno}{}
